% Generated draft: requires mathematical and source review.
\documentclass[11pt]{article}
\usepackage[margin=25mm]{geometry}
\usepackage{fontspec}
\setmainfont{Latin Modern Roman}
\usepackage{amsmath,amssymb,mathtools,unicode-math}
\setmathfont{Latin Modern Math}
\usepackage{enumitem,needspace,xurl,hyperref}
\setlength{\parindent}{0pt}
\setlength{\parskip}{5pt}
\setlength{\emergencystretch}{3em}
\newcommand{\N}{\mathbb{N}}
\newcommand{\Z}{\mathbb{Z}}
\newcommand{\Q}{\mathbb{Q}}
\newcommand{\R}{\mathbb{R}}
\newcommand{\C}{\mathbb{C}}
\newcommand{\F}{\mathbb{F}}
\newcommand{\A}{\mathbb{A}}
\newcommand{\PP}{\mathbb{P}}
\newcommand{\E}{\mathbb{E}}
\newcommand{\eps}{\varepsilon}
\newcommand{\dd}{\,\mathrm{d}}
\begin{document}
\section*{UnsolvedMath ID: 800726}
\textbf{Source problem code:} catalogue entry
\subsection{Definitions and mathematical statement}
Let \(F\) be an arbitrary field. Its Brauer group \(\operatorname{Br}(F)\) is the abelian group of Brauer equivalence classes of finite-dimensional central simple \(F\)-algebras, with addition induced by tensor product over \(F\). Here two such algebras \(A\) and \(B\) are Brauer equivalent if \(M_r(A)\cong M_s(B)\) as \(F\)-algebras for some positive integers \(r,s\). The identity is the class of \(F\).

For a prime \(p\), define the \(p\)-primary component by
\[
(\operatorname{Br}(F))_p
=\{\alpha\in\operatorname{Br}(F):p^n\alpha=0\text{ for some integer }n\geq 1\}.
\]
A subgroup \(D\) of an abelian group is divisible if, for every positive integer \(m\) and every \(d\in D\), there exists \(e\in D\) with \(me=d\). An elementary abelian \(2\)-group is an abelian group annihilated by \(2\), equivalently a vector space over \(\mathbb{F}_2\); no finiteness assumption is imposed.

\textbf{Conjecture 41 (attributed in the supplied source to Brumer--Rosen).} For every field \(F\) and every prime \(p\), at least one of the following holds:
\begin{enumerate}
\item \((\operatorname{Br}(F))_p=0\).
\item There exists a subgroup \(D\leq(\operatorname{Br}(F))_p\) such that \(D\neq 0\) and \(D\) is divisible.
\item \(p=2\) and \((\operatorname{Br}(F))_2\) is an elementary abelian \(2\)-group.
\end{enumerate}
The supplied statement imposes no restriction on the characteristic of \(F\). The alternatives are stated inclusively: the zero group also satisfies the definition of an elementary abelian \(2\)-group. The source's reference label \texttt{BrumerRosen} is supplied without its bibliographic entry, so the original Brumer--Rosen reference remains unresolved.
\subsection{Short English statement}
The conjecture asserts that, for any field and prime \(p\), the \(p\)-primary part of the Brauer group is zero or contains a nonzero divisible subgroup, with the additional possibility of an elementary abelian group when \(p=2\). The registry labels this conjecture open; that label is supplied source data and has not been independently verified.
\subsection{Sources}
\begin{itemize}
\item Local UnsolvedMath registry, permanent ID \(800726\), source problem code \texttt{catalogue entry}, titled \emph{Open Problems on Central Simple Algebras --- Conjecture 41}. The supplied record labels the problem \emph{open} and lists Auel, Brussel, Garibaldi, and Vishne with proposed year \(2010\).
\item Auel, Brussel, Garibaldi, and Vishne, \emph{Open Problems on Central Simple Algebras} (2010), Conjecture 41. Supplied source URL: \url{https://arxiv.org/abs/1006.3304}. The record reports access on \texttt{2026-10-16} and extraction by \texttt{arxiv\_latex}; these are supplied provenance details, not an independently verified source visit or literature assessment.
\item The supplied conjecture text attributes the statement to Brumer--Rosen using the reference label \texttt{BrumerRosen}. No full bibliographic entry for that reference is supplied, so the original attribution cannot be bibliographically resolved from this record alone.
\end{itemize}
\end{document}
